\chapter{附录：运行与验收说明}

\section{后端启动}

真实模型联调时，可使用如下命令启动后端：

\begin{lstlisting}
$env:MLLMPROJECT_USE_REAL_MODELS="1"
$env:MLLMPROJECT_QWEN3_MODEL_PATH="model"
python -m uvicorn mllmproject.api:app --host 127.0.0.1 --port 8000
\end{lstlisting}

无 GPU 或不需要真实模型时，可使用默认 mock backend 运行前后端流程。

\section{前端启动}

\begin{lstlisting}
cd frontend
npm.cmd install
npm.cmd run dev
\end{lstlisting}

浏览器访问 http://127.0.0.1:5173，即可使用上传、解析、知识库 chunk、PDF 预览和问答功能。

\section{评测命令}

\begin{lstlisting}
python scripts/run_benchmark_eval.py --mode text-rag
python scripts/run_benchmark_eval.py --mode mm-rag
\end{lstlisting}

评测结果保存到 data/eval/benchmarks/ 与 data/eval/results/。最终报告中的 20+20 结果来自仓库内已保存的评测输出。

\section{提交物}

本项目提交物包括：

\begin{itemize}
  \item 项目代码仓库与 README 运行说明。
  \item React + FastAPI Web Demo。
  \item 命令行评测脚本与实验结果。
  \item 开题报告、期末报告与展示 PPT。
  \item Demo 录屏材料。
\end{itemize}
